\chapter{ПРОЕКТИРОВАНИЕ АРХИТЕКТУРЫ РАСПРЕДЕЛЁННОГО ХРАНИЛИЩА}

В этой главе рассматривается архитектура отказоустойчивого распределённого хранилища «ключ--значение» (Key-Value Store), построенного на основе алгоритма консенсуса Raft. Описываются механизмы взаимодействия компонентов, алгоритмы маршрутизации запросов и методы обеспечения строгой консистентности (линеаризуемости) данных \cite{linearizability,raft_original}.

\section{Многоуровневая архитектура системы}

Разработанная система организована в виде слоистой архитектуры с четким разделением зон ответственности. Логически систему можно разделить на три функциональных уровня. 

Первый уровень составляет транспортная подсистема, отвечающая за вызовы удалённых процедур (RPC) между узлами кластера и обработку внешних HTTP-запросов от клиентов. Данный уровень инкапсулирует детали работы с сетью и сериализацией данных.
Второй уровень представляет собой ядро консенсуса. Оно изолированно реализует логику выбора лидера, репликации журнала, кворумной фиксации и поддержания согласованности истории операций, не опираясь на специфику прикладных данных. 
Третий уровень — хранилище "ключ--значение", выступающее в роли машины состояний, которую мы будем реплицировать.

Интеграция консенсусного ядра и машины состояний построена на принципах асинхронного конвейера. Алгоритм Raft принимает команды от клиентов, проводит их через процедуру репликации, и в момент достижения консенсуса передает зафиксированные команды в специализированный канал. Прикладной слой последовательно считывает эти команды и детерминированно применяет их к локальным данным. 

В качестве основного носителя данных прикладной машины состояний используется оперативная память. Данное архитектурное решение позволяет свести задержки при обработке прикладных запросов к минимуму. При этом фундаментальная гарантия долговечности данных не нарушается: вся входящая информация предварительно и синхронно сохраняется в энергонезависимый дисковый журнал на уровне ядра Raft. В случае аппаратного сбоя оперативное хранилище автоматически восстанавливает свое актуальное состояние путем повторного воспроизведения зафиксированных команд из дискового журнала \cite{kleppmann}.

\section{Механизмы изменения состояния}

Операции записи (такие как команда \texttt{Put}) инициируют изменение реплицированного состояния кластера. В соответствии с правилами алгоритма Raft, модификация данных может осуществляться исключительно через лидера. Поскольку клиент изначально может обратиться к любому случайному узлу (например, к последователю), в системе реализован механизм перенаправления. Если узел, получивший запрос на запись, не обладает лидерской ролью, он отклоняет запрос и возвращает клиенту сетевой адрес актуального координатора, что значительно сокращает время поиска точки входа в кластер.

При получении запроса лидером команда добавляется в журнал и рассылается остальным участникам. Критически важным архитектурным элементом является механизм синхронного ожидания: прикладной слой блокирует отправку ответа клиенту до тех пор, пока консенсус не подтвердит фиксацию записи на большинстве узлов. Это гарантирует, что клиент получит подтверждение об успешной операции только в том случае, если данные надежно сохранены в кластере и не могут быть утеряны при последующих сбоях.

В условиях асинхронной сети существует риск возникновения тайм-аута: команда записи может быть успешно зафиксирована кластером, но ответный сетевой пакет с подтверждением может быть утерян до доставки клиенту. Повторная отправка этой же команды клиентом может привести к её дублированному применению (например, к двойному списанию баланса).

Для устранения этой уязвимости архитектура хранилища реализует семантику строго однократного выполнения на базе идемпотентности клиентских сессий. К каждой операции модификации данных клиент прикрепляет уникальный идентификатор. Машина состояний кэширует результаты выполнения обработанных идентификаторов. При получении запроса с ранее встречавшимся идентификатором, прикладной слой распознает дубликат, блокирует повторное изменение данных и возвращает закэшированный результат предыдущего успешного выполнения \cite{kleppmann}.

\section{Обеспечение линеаризуемости при чтении данных}

Чтение значения напрямую из локальной памяти узла без предварительной проверки его статуса является небезопасным. Наиболее уязвимый сценарий возникает при нарушении связности сети. Лидер, оказавшийся в изолированном меньшинстве, может не знать о том, что в остальной части сети уже избран новый лидер и данные были изменены. Если изолированный узел ответит на запрос чтения из своей локальной оперативной памяти, клиент получит неверные, устаревшие данные, что нарушает свойство строгой консистентности (линеаризуемости).

Для обеспечения актуальности данных без выполнения дорогостоящих операций записи в дисковый журнал, в системе реализован механизм барьера чтения. Алгоритм чтения данных (\texttt{Get}) состоит из трех этапов \cite{raft_dissertation,linearizability}:

\begin{enumerate}
    \item Лидер фиксирует свой текущий индекс подтвержденных записей.
    \item Лидер рассылает параллельные пустые сигналы активности остальным узлам. Успешный сбор ответов от большинства доказывает, что лидер легитимен на текущий момент времени и не находится в сетевой изоляции.
    \item Лидер ожидает, пока прикладная машина состояний не применит все команды вплоть до зафиксированного ранее индекса, после чего возвращает данные из памяти.
\end{enumerate}

Особым краевым случаем является обработка чтений сразу после выбора нового координатора. Для инициализации актуального состояния и безопасного применения механизма барьера чтения, новый лидер обязан дождаться фиксации хотя бы одной (обычно пустой) операции в своей новой эпохе. В этот промежуток времени запросы на чтение блокируются алгоритмом до момента подтверждения консенсуса, что устраняет риск чтения данных по унаследованному, неактуальному индексу.

\section{Клиентский интерфейс}

Внешний интерфейс разработанного хранилища логически разделен на две независимые категории: операции с предметными данными (\texttt{Get}, \texttt{Put}) и административные операции управления составом кластера (\texttt{AddNode}, \texttt{RemoveNode}).

Все операции взаимодействия с кластером образуют предсказуемый синхронный контракт: клиент блокируется в ожидании ответа до тех пор, пока система не сможет гарантировать выполнение операции в соответствии с заявленными требованиями консистентности. Практически такая модель интерфейса близка к промышленным координационным сервисам и реестрам конфигураций, например etcd и ZooKeeper \cite{etcd_docs,zookeeper}. Для дискового прикладного слоя в дальнейшей эволюции системы могут применяться LSM-движки, такие как LevelDB и Pebble \cite{leveldb,pebble}.
